\scanpage{19}
% archive sheet number: 18
Soit $X$ un schéma algébrique lisse sur \deleted{\unsure{le}} un corps $k$ isomorphe à $\mathbb C$ \unsure{(}[de façon non précisée]. Alors
\[
 A^p(X)\xrightarrow{\unsure{P|\rho}_{\mathrm{alg}}}
 \underline{\underline H}^{2p}(X,\unsure{\Omega}_X^{\cdot})
\]
est bien défini. Le choix d'un isom $\sigma$ $k\xleftarrow{\sim}\mathbb C$ \deleted{\unsure{fixe}} permet de définir
\[
 P_\sigma^{(p)}:\quad
 A^p(X)\xrightarrow{\mathrm{can.}\,P^{\deleted{\unsure{\mathrm{an}}}}_{\mathrm{top}}\;\deleted{\unsure{P^{\mathrm{an}}}}}
 \underline{\underline H}^{2p}(X,\Omega_X^{\cdot})
\]
% slanted
\margin{$\mathrm{can}=$ De \uline{Rham} $+$ \uline{Gaga}}
qui se factorise par un sous-réseau $E_\sigma^{2i}=H^{2p}(X_\sigma,\mathbb Z)$ de $\underline{\underline H}^{2p}(X,\Omega_X^{\cdot})$. Or on a
\[
 P_\sigma^{(p)}=\left(\frac{1}{\sigma^{\deleted{\unsure{1|-1}}}(2i\pi)}\right)^p P_{\mathrm{alg}},
\]
ce qui montre que $P_\sigma^{(p)}$ dépend beaucoup de $\sigma$, \inserted{$\pi$ étant \uline{transcendant}} Dans le cas où $A^p(X)\xrightarrow{\substack{P_\sigma^{\unsure{(p)}}\\\sim}}H^{2p}(X_\sigma,\mathbb Z)$ est un isom (ex : $p=1$, $X=\mathbb P^1$) on voit de \unsure{$=$|même} que le réseau $E_\sigma^{2i}$, égal à $\frac{1}{\sigma(2i\pi)}P_{\mathrm{alg}}(A^p(X))$, dépend beaucoup de $\sigma$. \uline{NB} \unsure{Rempl.çant|Remplaçant} $\sigma$ par $\overline{\sigma}$ change $P_\sigma^{(p)}$ en $-P_\sigma^{(p)}$, donc ne change pas les $E_\sigma^i=H^{\deleted{\unsure{2i}}i}(X_\sigma,\mathbb Z)$

